%! The Victory of Orthogonality — Unit 7, Lesson 7.4 — Guided Notes
\documentclass[10pt]{article}
\usepackage{linalg-article}
\usepackage{linalg-boxes}
\usetikzlibrary{arrows.meta, positioning, calc}
\newcommand{\TallMath}[1]{$\displaystyle #1\rule[-1.4em]{0pt}{3.2em}$}
\newcommand{\uu}{\mathbf{u}}
\newcommand{\vv}{\mathbf{v}}
\newcommand{\xx}{\mathbf{x}}
\newcommand{\qq}{\mathbf{q}}
\newcommand{\zero}{\mathbf{0}}
\newcommand{\T}{^{\top}}

\begin{document}

\pageheader{Unit 7, Lesson 7.4}{Guided Notes}
\namedateperiod

\vspace{0.12in}

\begin{objectivebox}
\textbf{By the end of this lesson, I will be able to\dots}
\begin{itemize}\setlength{\itemsep}{2pt}
  \item recognize an orthogonal matrix by $Q\T Q=I$ and read off its free inverse $Q^{-1}=Q\T$;
  \item show that an orthogonal matrix preserves length, $|Q\xx|=|\xx|$;
  \item read the SVD $A=U\Sigma V\T$ as rotate--stretch--rotate, and explain why orthogonality is behind every result this unit.
\end{itemize}
\end{objectivebox}

\vspace{0.06in}

\begin{vocabbox}
\small
Fill in each term as we build it together.
\vspace{2pt}
\termblanklong{Orthogonal matrix (a square $Q$ with perpendicular unit columns --- $Q\T Q=I$)}
\termblanklong{Free inverse (for orthogonal $Q$, $Q^{-1}=Q\T$ --- just transpose)}
\termblanklong{Length-preserving ($|Q\xx|=|\xx|$ --- $Q$ moves vectors rigidly, no stretching)}
\end{vocabbox}

\begin{hookbox}
Look back at the whole unit: singular vectors came in \emph{perpendicular} pairs, compression and PCA both used
perpendicular directions, symmetric matrices had perpendicular eigenvectors.
\begin{itemize}\setlength{\itemsep}{1pt}
  \item What one idea has been doing the work behind every result this unit? \writeline
  \item Guess: what is special about a matrix whose columns are perpendicular unit vectors? \writeline
\end{itemize}
\end{hookbox}

\begin{notesbox}{1. Orthogonal Matrices: One Test, One Free Inverse}
An \textbf{orthogonal matrix} $Q$ is a square matrix whose columns are \textbf{perpendicular unit vectors}. All
of that is packed into a single equation --- form $Q\T Q$ (each entry is a dot product of two columns):
\[
  Q=\tfrac1{\sqrt2}\begin{bmatrix} 1 & 1\\ 1 & -1 \end{bmatrix}\ \Rightarrow\
  Q\T Q=\begin{bmatrix}\rule{0pt}{1.1em}\blank{0.8cm} & \blank{0.8cm}\\[2pt]\blank{0.8cm} & \blank{0.8cm}\end{bmatrix}=I.
\]
The diagonal $1$'s say each column has \textbf{length one}; the off-diagonal $0$'s say the columns are
\textbf{perpendicular}. The immediate payoff: since $Q\T Q=I$, the transpose \emph{is} the inverse,
\[
  Q^{-1}=\blank{1.4cm}.
\]
No elimination, no determinant --- the inverse is \textbf{free}.
\end{notesbox}

\begin{notesbox}{2. Orthogonal Matrices Preserve Length}
Watch what $Q$ does to the length of any vector. Using $Q\T Q=I$:
\[
  |Q\xx|^2=(Q\xx)\T(Q\xx)=\xx\T Q\T Q\,\xx=\xx\T\xx=|\xx|^2.
\]
So $|Q\xx|=|\xx|$ for \emph{every} $\xx$: an orthogonal matrix moves vectors \textbf{rigidly} --- a rotation or a
reflection --- with no stretching and no distortion. Check it on $\xx=(3,1)$ (length $\sqrt{10}$):
\[
  Q\xx=\tfrac1{\sqrt2}\begin{bmatrix}\rule{0pt}{1.1em}\blank{0.7cm}\\[2pt]\blank{0.7cm}\end{bmatrix},\qquad
  |Q\xx|^2=\blank{1.0cm}=|\xx|^2.
\]
\end{notesbox}

\begin{notesbox}{3. The SVD is Rotate--Stretch--Rotate}
\begin{minipage}[t]{0.52\linewidth}\vspace{0pt}
Every matrix factors as $A=U\Sigma V\T$ with $U,V$ \textbf{orthogonal} and $\Sigma$ \textbf{diagonal}. Reading
right to left, applying $A$ to $\xx$ is three moves:
\begin{itemize}[leftmargin=1.2em, nosep]
  \item $V\T$ --- \textbf{turns} $\xx$ to line up with the axes (orthogonal: no stretch);
  \item $\Sigma$ --- \textbf{stretches} along the axes by $\sigma_1,\sigma_2$;
  \item $U$ --- \textbf{turns} to the output direction (orthogonal: no stretch).
\end{itemize}
\vspace{2pt}
\textbf{All the stretching lives in $\Sigma$}; the orthogonal factors only rotate. Our spine
$A=\begin{bsmallmatrix} 2 & 1\\ 1 & 2 \end{bsmallmatrix}$ from Unit~6 is exactly
$Q\,\Sigma\,Q\T$ with $\Sigma=\mathrm{diag}(3,1)$: rotate, stretch by $3$ and $1$, rotate back.
\end{minipage}\hfill
\begin{minipage}[t]{0.44\linewidth}\vspace{0pt}
\centering
\begin{tikzpicture}[font=\scriptsize, node distance=6mm,
    box/.style={draw=burgundy, fill=blush, rounded corners=1pt, minimum width=9mm, minimum height=6.5mm, align=center},
    ar/.style={-{Latex[length=1.8mm]}, burgundy, thick}]
  \node[box] (x)  {$\xx$};
  \node[box, right=of x] (v)  {$V\T$};
  \node[box, right=of v] (s)  {$\Sigma$};
  \node[box, right=of s] (u)  {$U$};
  \node[box, right=of u] (ax) {$A\xx$};
  \draw[ar] (x)--(v);
  \draw[ar] (v)--(s);
  \draw[ar] (s)--(u);
  \draw[ar] (u)--(ax);
  \node[burgundy] at ($(v)+(0,0.85)$) {turn};
  \node[burgundy] at ($(s)+(0,0.85)$) {stretch};
  \node[burgundy] at ($(u)+(0,0.85)$) {turn};
\end{tikzpicture}\\[3pt]
{\scriptsize $V\T$ and $U$ only rotate; $\Sigma$ does all the stretching.}
\end{minipage}
\end{notesbox}

\begin{notesbox}{4. The Victory --- and Why It Matters}
\textbf{Why orthogonal matrices win.} Because $U$ and $V$ are orthogonal, they are \emph{reversible for free}
($Q^{-1}=Q\T$), they \emph{never change length}, and --- crucially for computing --- they \emph{never amplify
errors}. That last fact is the ``revolution'': modern algorithms deliberately compute with orthogonal matrices
because rounding mistakes stay small instead of blowing up.

\vspace{2pt}
\textbf{The thread through the course.} Orthogonality has been the quiet engine all along: perpendicular
projections and least squares (Unit~4), perpendicular eigenvectors of symmetric matrices (Unit~6), and now
perpendicular singular vectors for \emph{every} matrix, any shape (Unit~7).

\vspace{2pt}
\textbf{The victory.} The SVD $A=U\Sigma V\T$ says every matrix is really just \textbf{rotate--stretch--rotate}.
The two rotations are orthogonal --- perfect, distortion-free --- and one diagonal $\Sigma$ carries all the
stretching. That is why the SVD, image compression, and PCA all work: they run on orthogonality.
\end{notesbox}

\begin{practicebox}
Show your work.
\begin{enumerate}[leftmargin=1.6em, itemsep=4pt]
  \item Is $Q=\tfrac1{\sqrt2}\begin{bsmallmatrix} 1 & -1\\ 1 & 1 \end{bsmallmatrix}$ orthogonal? Compute $Q\T Q$ to decide, and write $Q^{-1}$. \writeline
  \item An orthogonal matrix $Q$ acts on a vector with $|\xx|=5$. What is $|Q\xx|$, and why? \writeline
  \item In $A=U\Sigma V\T$, which factor does the stretching, and which factors only rotate? \writeline
\end{enumerate}
\end{practicebox}

\end{document}
